
\section{Evaluations}
\label{sec:exp}

% We present three main empirical findings with \method.
% First, we qualitatively and quantitatively compare \method to existing benchmarks, and validate that \method fixes the identified biases and it is more faithful in evaluating a model's vision-language compositionality (Sec.~\ref{sec:exp-bias}).
% Second, we re-evaluate recent methods proposed to improve compositionality of VLMs, and show that their effectiveness is overestimated on existing benchmarks (Sec.~\ref{sec:exp-reeval}).
% Lastly, we provide a comprehensive evaluation on a wide array of pretrained CLIP models, and reveal that current models still lack compositionality (Sec.~\ref{sec:exp-eval}).

In this section, we qualitatively and quantitatively compare \method to existing benchmarks (Sec.~\ref{sec:exp-bias}), 
re-evaluate recent methods proposed to improve compositionality of vision-language models (Sec.~\ref{sec:exp-reeval}), and comprehensively evaluate a wide array of pretrained CLIP models (Sec.~\ref{sec:exp-eval}).

% \textbf{Setup.} Existing compositionality benchmarks are generated from different source datasets~\cite{krishna2017visual,lin2014microsoft}, and thus contain different original images and positive texts. In addition, the exact formulations of the evaluation retrieval tasks differ slightly between existing benchmarks.
% To systematically and fairly compare \method with \emph{all} existing benchmarks, we normalize the benchmarks in order to conduct careful controlled experiments.

% \textbf{Normalizing existing benchmarks with \hncoco.}
% We normalize existing benchmarks by reproducing their data generation workflow using COCO~\cite{lin2014microsoft} dataset as the underlying base image-text dataset, as in \method. We refer to this reproduced hard negative dataset based on COCO as \hncoco. In Appendix~\ref{app:bms}, we show the mapping from the existing benchmarks to their corresponding reproduced versions in \hncoco. In addition, we standardize the evaluation task as retrieving the correct caption from \emph{two} possible choices, i.e, a positive text and a hard negative.
% This normalization sets the positive texts fixed for all benchmarks, including \method, and enables us to directly compare the quality of the hard negatives generated in different benchmarks.

To systematically and fairly compare \method with existing benchmarks, we normalize the benchmarks by reproducing their data generation workflow using COCO~\cite{lin2014microsoft} as in \method. We utilize source code from CREPE~\cite{ma2022crepe} to generate \replace, \swap, \negate\ hard negatives and take \shuffle\ hard negatives released in ARO~\cite{yuksekgonul2023when}. We refer to this reproduced dataset as \hncoco.
% In Appendix~\ref{app:bms}, we show the mapping from the existing benchmarks to their corresponding reproduced versions in \hncoco.
In addition, we standardize the evaluation task as retrieving the correct caption from \emph{two} possible choices, \ie, a positive text and a hard negative.
This normalization sets the positive texts fixed for all benchmarks, including \method.

\subsection{\method significantly reduces dataset biases}
\label{sec:exp-bias}

\begin{table*}[!t]
  \centering
  \small
  \caption{We present example positive texts and their hard negatives in \hncoco\ (generated using existing procedures) and \method (generated with ChatGPT). \method brings significant improvements in commonsense and fluency.}
  \resizebox{\linewidth}{!}{
    \begin{tabular}{l lll} 
    \toprule
    Hard-Negative Type & Text Type & Commonsense & Fluency\\
    \midrule
    
     & Original & Two adult bears play fight in the water. & A man sitting in front of a laptop computer.\\
     \replace & \hncoco & Two adult bears play fight in the soda. & A man sitting around front of a laptop computer.\\
     & \method & A flock of ducks play fight in the water. & A man standing in front of a laptop computer.\\

     \midrule
    
    & Original & A woman standing behind a fence looking at an elephant. & Man swinging tennis racket while group of people watches. \\
     \swap & \hncoco & A fence standing behind a woman looking at an elephant. & Group swinging tennis racket while man of people watches. \\
     & \method & An elephant standing behind a fence looking at a woman. & Group of people swinging tennis racket while man watches. \\
    
    
    \midrule
    
    & Original & A teddy bear next to a stuffed fish. & A red fire hydrant on a city sidewalk. \\
     \negate\ / \add & \hncoco &  A teddy bear next to a stuffed fish. There is no teddy bear. & A red fire not hydrant on a city sidewalk.\\
     & \method & A teddy bear and a stuffed fish and a robot toy. & A red fire hydrant and a trash can on a city sidewalk. \\

\bottomrule 
\end{tabular}
}
  
  \label{tab:sugar_crepe_examples}
\end{table*}


\textbf{\method generates more sensical and fluent hard negatives.}
We validate that \method generates higher quality hard negative texts by leveraging ChatGPT than previous rule-based approaches.
Qualitatively, in Table~\ref{tab:sugar_crepe_examples}, we observe that the hard negatives in \method are more sensical and fluent compared to hard negatives in \hncoco.
We report human evaluation results in Appendix~\ref{app:sugar} that show on an average of $35\%$ of examples, hard negatives in \method have \emph{strictly} higher quality than \hncoco\ in terms of commonsense and fluency. For instance, on \swap, humans judge that \method wins $68\%$ over \hncoco\ and ties on $28\%$ of examples in terms of commonsense.
Quantitatively, in Table~\ref{tab:sugar_crepe_eval_quant_auto}, we compare the commonsense and grammar scores averaged over the hard negative texts in both \hncoco\ and \method. We see \method has much higher average scores than \hncoco. Additionally, pairwise comparisons show that \method has higher commonsense and grammar scores than \hncoco\ on $86\%$ of examples on average.


\begin{table*}[!t]
  \centering
  \small
  \caption{We compare the commonsense and grammar scores on hard negatives in \hncoco\ and \method. We report both their respective average scores and the ratio where \method has higher score than \hncoco\ in pairwise comparison. Overall, \method has hard negatives with better commonsense and grammar.}
  \scalebox{0.8}{
    \begin{tabular}{ll ccc} 
    \toprule
    & & \multicolumn{2}{c}{Average Score} \\
    \cmidrule{3-4} 
    Hard-negative Type & Metric & \hncoco & \method & Pairwise Better Ratio \\
    \midrule
    
    \multirow{2}{*}{\replace} & Commonsense & 37.46 & 50.21 & 77.71 \\
    & Grammar & 76.79 & 88.96 & 86.85\\
    
    
    \midrule
    
    \multirow{2}{*}{\swap} & Commonsense  & 23.09 & 41.57 & 78.76\\
    & Grammar & 45.67 & 80.46 & 87.02\\
    
    
    \midrule
    
    \multirow{2}{*}{\negate\ / \add} & Commonsense &  25.24 & 50.20 & 87.24\\
    & Grammar & 65.09 & 90.07 & 95.03\\
    
    
\bottomrule 

\end{tabular}
}
  
  \label{tab:sugar_crepe_eval_quant_auto}

\end{table*}


\textbf{\method disentangles the identified exploitable biases.}
% While the hard negatives generated from ChatGPT in \method already much reduces the biases compared to existing benchmarks
We show that the final \method evaluation set maximally reduces the identified biases that could be exploited undesirably to achieve improvements on a benchmark.
Figure~\ref{fig:vera_gap_ours} visualizes the Vera/Grammar score gap distributions.
We compare the distributions between \hncoco\ and \method (before and after adversarial refinement).
First, We see that by leveraging ChatGPT, the hard negative texts in \method already have lower biases than \hncoco\ before adversarial refinement, \ie, the score gap distribution is more centered around zero.
Furthermore, we see that after adversarial refinement, the score gap distributions on the final \method evaluation set are symmetric around zero. This implies that the previously identified artifacts can no longer be exploited to infer the positive texts.
As a result, we show that the previous commonsense and grammar attacks that are extremely successful on existing benchmarks do not work on \method. As shown in Table~\ref{tab:eval_ours}, these blind models now consistently rank the \emph{last} on \method as compared to other pretrained CLIP models.


\begin{figure}[!t]
    \centering
    \includegraphics[width=\linewidth]{imgs/score_gap_ours.png}
    \caption{We compare the Vera (top row) and Grammar (bottom row) score gap distributions between \hncoco\ (leftmost column), \method without adversarial refinement (middle), and \method (rightmost). Top row: We see that Vera score gap distribution shifts from the positive spectrum to more centered around zero from \hncoco\ to \method without refinement. After adversarial refinement, we ensure the score gap distribution is centered around zero on \method. Bottom row: Similarly, from \hncoco\ to \method, we see the Grammar score gap distribution shifts from the positive spectrum to centered around zero.}
    \label{fig:vera_gap_ours}
\end{figure}

% \begin{figure}[!t]
%     \centering
%     \includegraphics[width=\linewidth]{imgs/vera_gap_ours.png}
%     \caption{Vera score: \hncoco\ v.s. \method candidate v.s. \method. \cy{only keep high-level types.} \aniline{Instead of canddiate, can we call it With No AF ?}}
%     \label{fig:vera_gap_ours}
% \end{figure}

% \begin{figure}[!t]
%     \centering
%     \includegraphics[width=\linewidth]{imgs/grammar_gap_ours.png}
%     \caption{Grammar score: \hncoco\ v.s. \method candidate v.s. \method \rk{Make sure that the text is large enough in these figures. And also make sure that each caption lays out the make takeaway a reader should have when viewing each figure.}}
%     \label{fig:grammar_gap_ours}
% \end{figure}



\subsection{Re-evaluating recent methods for improving compositionality}
\label{sec:exp-reeval}
Given the vulnerability of existing compositionality benchmarks, it is unclear whether recently proposed methods that show state-of-the-art performances on these benchmarks are indeed effective. Thus, we re-evaluate these methods with \method.

\textbf{Hard negative augmented training.} 
Specifically, we focus on evaluating one common \emph{data-augmentation} approach considered in~\cite{yuksekgonul2023when,doveh2023teaching}, where the core idea is to explicitly create hard negatives and train the model to distinguish them. We broadly refer to this training scheme as \negclip following \cite{yuksekgonul2023when}.
We evaluate two \negclip training schemes: finetuning and training from scratch.
For finetuning, in addition to taking the model released in \cite{yuksekgonul2023when}, we finetune another three \negclip\ models (using ViT-B/32 following~\cite{yuksekgonul2023when}) with three respective types of hard negatives (\ie, \replace, \swap, \negate) generated using CREPE's~\cite{ma2022crepe} source code.
For training from scratch, we use RN50 as the base model and train variants of \negclip by augmenting the training examples with different types of hard negatives.
We perform both training and finetuning on COCO~\cite{lin2014microsoft}.

\textbf{Improvements are overestimated due to unintentionally overfitting.}
In Table~\ref{tab:re_eval}, we first see that \negclip\ finetuned models show significant improvements on \hncoco, boosting the performance more than 10\% compared to standard CLIP finetuning on 11 out of 16 cases (highlighted in green). The lifts are especially large when the hard negative type used in finetuning matches that used in evaluation, where \negclip\ finetuned models can achieve near human-level performances. For instance, by finetuning with \replace\ hard negatives, \negclip\ reaches 94\% on \hncoco\ evaluated with \replace\ hard negatives (human performance is 95\%).
While the results on \hncoco\ suggest that \negclip\ is seemingly sufficient in equipping models with strong compositionality, we however see that the improvements brought by \negclip\ are much smaller on \method. In fact, none of the improvements on \method is larger than 10\%, and the best performing \negclip\ finetuned models still have large gaps to human-level performances, \eg, best \negclip\ model lags behind human by 23\% on \method's \swap\ hard negatives.
Similarly, when trained from scratch, we observe the same trend that \negclip's improvements are much larger on \hncoco\ than on \method. The improvements on \hncoco\ are again most pronounced when the training and testing hard negative type matches.

We attribute the stark contrast in \negclip's effectiveness on \hncoco\ and \method to model's unintentional overfitting: The \negclip\ models learned to exploit artifacts that can be used to easily distinguish hard negatives from positives on \hncoco, instead of actually improving compositionality.
Thus, when evaluated on \method where the artifacts are removed, the improvement from \negclip\ drastically reduces.
These results imply that \negclip's effectiveness is overestimated on existing benchmarks, and we may still need further innovations to fundamentally improve a model's compositionality.
% Finally, by comparing a model's performance on \hncoco\ and \method, we see that the exaggerated improvements resulting from overfittings on \hncoco\ are much reduced on \method, as \method removes the exploitable biases.


\begin{table*}[h]
  \centering
  \small
  \caption{Re-evaluating hard negative augmented training shows that the method's improvements on existing benchmarks (\hncoco) are hugely overestimated, particularly when the test hard negative type matches the one used in training, which can be attributed to overfitting the artifacts.\\Color notations: \colorbox{tablegreen}{Gains compared to standard CLIP (finetuned / from scratch) $>$ 10\%}.}
  \begin{adjustbox}{max width=\textwidth}
    \begin{tabular}{lll cccc ccc} 
    \toprule
     & & & \multicolumn{4}{c}{\bf \hncoco} &  \multicolumn{3}{c}{\bf \method} \\ 
     
     \cmidrule(lr){4-7} \cmidrule(lr){8-10}
     
    \textbf{Model} & \textbf{Training} & \textbf{Hard Negative Used} & \textbf{\replace} & \textbf{\swap} & \textbf{\negate} & \textbf{\shuffle} &  \textbf{\replace} & \textbf{\swap} & \textbf{\add} \\ 
    \midrule

Human & & & 95.33 & 100 & 99.33 & 96.00 & 98.67 & 99.50 & 99.00 \\ \midrule


\multirow{5}{*}{ViT-B/32} & Pretrained & N/A & 75.71 & 71.58 & 76.89 & 72.06 & 80.76 & 63.27 & 75.09\\

& CLIP finetuned & N/A & 77.06 & 68.81 & 61.19 & 63.04 & 84.76 & 70.83 & 85.58\\
\cmidrule{2-10}


& \multirow{4}{*}{\negclip\ finetuned} & \replace & \colorbox{tablegreen}{94.51} & \colorbox{tablegreen}{90.04} & \colorbox{tablegreen}{85.06} & \colorbox{tablegreen}{88.15} & 88.27 & 74.89 & 90.16 \\
& & \swap & 82.88 & \colorbox{tablegreen}{94.48} & \colorbox{tablegreen}{77.57} & \colorbox{tablegreen}{87.00} & 85.54 & 76.21 & 86.56\\
& & \negate & 77.24 & 68.91 & \colorbox{tablegreen}{99.54} & 64.28 & 84.97 & 70.29 & 85.84 \\
&  & Released in~\cite{yuksekgonul2023when} & 85.72 & \colorbox{tablegreen}{94.35} & \colorbox{tablegreen}{83.51} & \colorbox{tablegreen}{90.45} & 85.36 & 75.33 & 87.29 \\ 

\midrule
\multirow{5}{*}{RN50}
& CLIP from scratch & N/A & 69.93 & 59.96 & 55.36 & 68.78 & 69.54 & 60.33 & 67.63\\
\cmidrule{2-10}
& \multirow{4}{*}{\negclip\ from scratch} & \replace & \colorbox{tablegreen}{89.04} & 66.51 & 60.90 & 75.23 & 74.32 & 62.65 & 72.92\\
& & \swap & 72.33 & \colorbox{tablegreen}{92.29} & 64.51 & \colorbox{tablegreen}{84.84} & 73.31 & 68.35 & 71.93\\
& & \negate & 70.09 & 60.29 & \colorbox{tablegreen}{99.45} & 69.03 & 72.74 & 60.89 & 70.47\\
& & \textsc{Rep} + \textsc{Sw} + \textsc{Neg} &  \colorbox{tablegreen}{86.30} & \colorbox{tablegreen}{88.60} & \colorbox{tablegreen}{99.34} & \colorbox{tablegreen}{82.93} & 75.26 & 67.69 & 73.08\\

\bottomrule 
\end{tabular}
\end{adjustbox}
  
  \label{tab:re_eval}

\end{table*}


% \begin{table*}[h]
%   \centering
%   \small
%   \caption{Re-evaluating hard-negative augmented training shows that the method's improvements on existing benchmark (\hncoco) are hugely overestimated, particularly when the test hard negative type matches the one used in training, which can be attributed to overfitting the artifacts. \\Color notations: \colorbox{tablegreen}{Gains compared to standard CLIP (finetuned / from scratch) $>$ 10$\%$}.}
%   \begin{adjustbox}{max width=\textwidth}
%     \begin{tabular}{lll cccc ccc} 
%     \toprule
%      & & & \multicolumn{4}{c}{\bf \hncoco} &  \multicolumn{3}{c}{\bf \method} \\ 

     
%      \cmidrule(lr){4-7} \cmidrule(lr){8-10}
     
%     \textbf{Model} & \textbf{Training} & \textbf{Hard Negative Used} & \textbf{\shuffle}~\cite{yuksekgonul2023when} & \textbf{\replace} & \textbf{\swap} & \textbf{\negate} &  \textbf{\replace} & \textbf{\swap} & \textbf{\add} \\ 
%     \midrule

% Human & & & 96.00 & 95.33 & 100 & 99.33 & 98.67 & 99.50 & 99.00 \\ \midrule


% \multirow{5}{*}{ViT-B/32} & Pretrained & N/A & 72.06 & 75.71 & 71.58 & 76.89 & 80.76 & 63.27 & 75.09\\

% & CLIP finetuned & N/A & 63.04 & 77.06 & 68.81 & 61.19 & 84.76 & 70.83 & 85.58\\
% \cmidrule{2-10}

% & \multirow{4}{*}{\negclip\ finetuned} & Released in~\cite{yuksekgonul2023when} & \colorbox{tablegreen}{90.45} & 85.72 & \colorbox{tablegreen}{94.35} & \colorbox{tablegreen}{83.51} & 85.36 & 75.33 & 87.29 \\ 
% &  & \replace & \colorbox{tablegreen}{88.15} & \colorbox{tablegreen}{94.51} & \colorbox{tablegreen}{90.04} & \colorbox{tablegreen}{85.06} & 88.27 & 74.89 & 90.16 \\
% & & \swap & \colorbox{tablegreen}{87.00} & 82.88 & \colorbox{tablegreen}{94.48} & \colorbox{tablegreen}{77.57} & 85.54 & 76.21 & 86.56\\
% & & \negate & 64.28 & 77.24 & 68.91 & \colorbox{tablegreen}{99.54} & 84.97 & 70.29 & 85.84 \\
%  % \cmidrule{2-10}


% \midrule
% \multirow{5}{*}{RN50}
% & CLIP from scratch & N/A & 68.78 & 69.93 & 59.96 & 55.36 & 69.54 & 60.33 & 67.63\\
% \cmidrule{2-10}
% & \multirow{4}{*}{\negclip\ from scratch} & \replace & 75.23 & \colorbox{tablegreen}{89.04} & 66.51 & 60.90 & 74.32 & 62.65 & 72.92\\
% & & \swap & \colorbox{tablegreen}{84.84} & 72.33 & \colorbox{tablegreen}{92.29} & 64.51 & 73.31 & 68.35 & 71.93\\
% & & \negate & 69.03 & 70.09 & 60.29 & \colorbox{tablegreen}{99.45} & 72.74 & 60.89 & 70.47\\
% & & \textsc{Rep} + \textsc{Sw} + \textsc{Neg} & \colorbox{tablegreen}{82.93} & \colorbox{tablegreen}{86.30} & \colorbox{tablegreen}{88.60} & \colorbox{tablegreen}{99.34} & 75.26 & 67.69 & 73.08\\


% \bottomrule 
% \end{tabular}
% \end{adjustbox}
  
%   \label{tab:re_eval}

% \end{table*}



\subsection{Comprehensive evaluations on existing pretrained vision-language models}
\label{sec:exp-eval}
% \begin{table*}[t]
%   \centering
%   \small
%   \caption{Our evaluation of pretrained vision-language models on \method shows that they demonstrate compositionality on some hard negatives but are far from human performance on others, especially on hard negatives of \swap\ form or ones perturbing attributes and relations.}
%   \scalebox{0.6}{
%     \begin{tabular}{llrr ccccccc} 
%     \toprule
%      &&&& \multicolumn{2}{c}{\bf \add} &  \multicolumn{2}{c}{\bf \swap}&\multicolumn{3}{c}{\bf \replace}\\ 
%      \cmidrule(lr){5-6}\cmidrule(lr){7-8}\cmidrule(lr){9-11}
     
%     \textbf{Source} &\textbf{Model} &\textbf{Data Size} &\textbf{Model Size (M)} & \textbf{Object} & \textbf{Attribute} & \textbf{Object} & \textbf{Attribute} & \textbf{Object} & \textbf{Attribute} & \textbf{Relation}  \\ 
%     \midrule

% \multicolumn{2}{c}{Human} & & & 91 & 80 & 90 & 93 & 95 & 92 & 90 \\

% \midrule\midrule

% \multirow{2}{*}{Text-only model} &  Vera~\cite{liu2023vera} & & & 49.30 & 49.49 & 49.17 & 49.57 & 49.29 & 49.74 & 49.38 \\

% & Grammar~\cite{morris2020textattack} & & & 50.00 & 50.00 & 50.00 & 50.00 & 50.00 & 50.00 & 50.00 \\

% \midrule\midrule

% \multirow{7}{*}{OpenAI~\cite{RadfordKHRGASAM21}}    
% & RN50& \multirow{7}{*}{400M} & 102& 71.27& 60.97& 65.62& 66.41& 88.13& 76.49& 64.92\\
% & RN101& & 120& 73.15& 58.92& 60.55& 62.05& 89.10& 78.96& 62.54\\
% & ViT-B-32& & 151& 74.24& 61.85& 63.56& 63.34& 87.32& 75.13& 64.83\\
% & RN50x4& & 178& 76.23& 61.71& 61.78& 64.12& 90.01& 77.17& 63.95\\
% & RN50x16& & 291& 77.56& 64.13& 64.38& 64.41& 90.87& 77.09& 64.44\\
% & ViT-L-14& & 428& 75.49& 64.64& 60.27& 63.34& 90.82& 75.55& 62.94\\
% & RN50x64& & 623& 79.90& 66.17& 61.92& 63.48& 91.73& 78.02& 66.37\\


% \midrule

% \multirow{6}{*}{LAION~\cite{laion5b}}   
% & roberta-ViT-B-32& \multirow{4}{*}{2B}& 212& 84.35& 72.21& 64.52& 66.12& 89.25& 80.75& 67.87\\
% & ViT-H-14&  & 986& 88.84& 75.56& 66.99& 66.83& 93.05& 82.11& 67.61\\
% & ViT-g-14& & 1367& 88.60& 77.09& 63.15& 65.76& 92.75& 82.20& 68.31\\
% & ViT-bigG-14& & 2540& 88.99& 74.81& 66.58& 67.83& 93.86& 84.41& 70.03\\ \cline{3-3}
% & xlm-roberta-base-ViT-B-32& \multirow{2}{*}{5B}& 366& 84.35& 73.09& 63.56& 65.91& 89.96& 79.90& 65.67\\
% & xlm-roberta-large-ViT-H-14& & 1193& 90.83& 76.12& 65.48& 66.48& 93.86& 82.37& 68.27\\


% \midrule

% \multirow{4}{*}{DataComp~\cite{gadre2023datacomp}}   
% & \texttt{small:}ViT-B-32& 13M& 151& 54.02& 57.25& 51.51& 52.00& 57.10& 55.79& 53.79\\
% & \texttt{medium:}ViT-B-32& 128M& 151 &65.61& 61.94& 57.53& 55.21& 73.53& 66.52& 54.97\\
% & \texttt{large:}ViT-B-16& 1B& 150& 81.26& 69.98& 58.22& 58.56& 88.79& 75.21& 61.18\\
% & \texttt{xlarge:}ViT-L-14& 13B& 428& 88.10& 76.02& 63.97& 64.48& 92.95& 79.81& 66.68\\

% \bottomrule 
% \end{tabular}
% }
  
%   \label{tab:eval_ours}

% \end{table*}


% \begin{table*}[t]
%   \centering
%   \small
%   \caption{\cy{need update.}Our evaluation of pretrained CLIP models on \method shows that they demonstrate compositionality on some hard negatives but are far from human performance on others, especially on \swap\ hard negatives or ones perturbing attributes and relations (also illustrated in Figure~\ref{fig:imagenet_correlation}: lower overall performance on \swap, and lower performances on attributes/relations compared to objects).}
%   \scalebox{0.6}{
%     \begin{tabular}{llrr ccccccc} 
%     \toprule
%      &&&& \multicolumn{3}{c}{\bf \replace}  &  \multicolumn{2}{c}{\bf \swap} & \multicolumn{2}{c}{\bf \add} \\ 
%      \cmidrule(lr){5-7}\cmidrule(lr){8-9}\cmidrule(lr){10-11}
     
%     \textbf{Source} &\textbf{Model} &\textbf{Data Size} &\textbf{Model Size (M)} & \textbf{Object} & \textbf{Attribute} & \textbf{Relation} & \textbf{Object} & \textbf{Attribute} & \textbf{Object} & \textbf{Attribute}   \\ 
%     \midrule

% \multicolumn{2}{c}{Human} & & & 95 & 92 & 90  & 90 & 93 & 91 & 80\\

% \midrule\midrule

% \multirow{2}{*}{Text-only model} &  Vera~\cite{liu2023vera} & & & 49.29 & 49.74 & 49.38 & 49.17 & 49.57 &  49.30 & 49.49\\

% & Grammar~\cite{morris2020textattack} & & & 50.00 & 50.00 & 50.00 & 50.00 & 50.00 & 50.00 & 50.00 \\

% \midrule\midrule

% \multirow{7}{*}{OpenAI~\cite{RadfordKHRGASAM21}}    
% & RN50& \multirow{7}{*}{400M} & 102& 88.13& 76.49& 64.92 & 65.62& 66.41& 71.27& 60.97\\
% & RN101& & 120& 89.10& 78.96& 62.54 & 60.55& 62.05& 73.15& 58.92\\
% & ViT-B-32& & 151& 87.32& 75.13& 64.83 & 63.56& 63.34& 74.24& 61.85\\
% & RN50x4& & 178& 90.01& 77.17& 63.95 & 61.78& 64.12& 76.23& 61.71\\
% & RN50x16& & 291& 90.87& 77.09& 64.44 & 64.38& 64.41& 77.56& 64.13\\
% & ViT-L-14& & 428& 90.82& 75.55& 62.94 & 60.27& 63.34& 75.49& 64.64\\
% & RN50x64& & 623& 91.73& 78.02& 66.37 & 61.92& 63.48& 79.90& 66.17\\


% \midrule

% \multirow{6}{*}{LAION~\cite{laion5b}}   
% & roberta-ViT-B-32& \multirow{4}{*}{2B}& 212& 89.25& 80.75& 67.87 & 64.52& 66.12& 84.35& 72.21\\
% & ViT-H-14&  & 986& 93.05& 82.11& 67.61 & 66.99& 66.83& 88.84& 75.56\\
% & ViT-g-14& & 1367& 92.75& 82.20& 68.31 & 63.15& 65.76& 88.60& 77.09\\
% & ViT-bigG-14& & 2540& 93.86& 84.41& 70.03 & 66.58& 67.83& 88.99& 74.81\\ \cline{3-3}
% & xlm-roberta-base-ViT-B-32& \multirow{2}{*}{5B}& 366& 89.96& 79.90& 65.67 & 63.56& 65.91& 84.35& 73.09\\
% & xlm-roberta-large-ViT-H-14& & 1193& 93.86& 82.37& 68.27 & 65.48& 66.48& 90.83& 76.12\\


% \midrule

% \multirow{4}{*}{DataComp~\cite{gadre2023datacomp}}   
% & \texttt{small:}ViT-B-32& 13M& 151& 57.10& 55.79& 53.79 & 51.51& 52.00& 54.02& 57.25\\
% & \texttt{medium:}ViT-B-32& 128M& 151 & 73.53& 66.52& 54.97 & 57.53& 55.21& 65.61& 61.94\\
% & \texttt{large:}ViT-B-16& 1B& 150& 88.79& 75.21& 61.18 & 58.22& 58.56& 81.26& 69.98\\
% & \texttt{xlarge:}ViT-L-14& 13B& 428& 92.95& 79.81& 66.68 & 63.97& 64.48& 88.10& 76.02\\

% \bottomrule 
% \end{tabular}
% }
  
%   \label{tab:eval_ours}

% \end{table*}


\begin{table*}[t]
  \centering
  \small
  \caption{Our evaluation of pretrained CLIP models on \method shows that they demonstrate compositionality on some hard negatives but are far from human performance on others, especially on \swap\ hard negatives or ones perturbing attributes and relations (also illustrated in Figure~\ref{fig:imagenet_correlation}: lower overall performance on \swap, and lower performances on attributes/relations compared to objects).}
  \scalebox{0.6}{
    \begin{tabular}{llrr ccccccc} 
    \toprule
     &&&& \multicolumn{3}{c}{\bf \replace}  &  \multicolumn{2}{c}{\bf \swap} & \multicolumn{2}{c}{\bf \add} \\ 
     \cmidrule(lr){5-7}\cmidrule(lr){8-9}\cmidrule(lr){10-11}
     
    \textbf{Source} &\textbf{Model} &\textbf{Data Size} &\textbf{Model Size (M)} & \textbf{Object} & \textbf{Attribute} & \textbf{Relation} & \textbf{Object} & \textbf{Attribute} & \textbf{Object} & \textbf{Attribute}   \\ 
    \midrule

\multicolumn{2}{c}{Human} & & & 100 & 99 & 97  & 99 & 100 & 99 & 99\\

\midrule\midrule

\multirow{2}{*}{Text-only model} &  Vera~\cite{liu2023vera} & & & 49.39 & 49.62 & 49.36 & 49.19 & 49.40 &  49.42 & 49.57\\

& Grammar~\cite{morris2020textattack} & & & 50.00 & 50.00 & 50.00 & 50.00 & 50.00 & 50.00 & 50.00 \\

\midrule\midrule

\multirow{7}{*}{OpenAI~\cite{RadfordKHRGASAM21}}    
& RN50& \multirow{7}{*}{400M} & 102& 91.77& 80.58& 69.99& 61.79& 68.47& 74.54& 69.65\\
& RN101& & 120& 92.49& 83.88& 67.07& 56.50& 65.92& 75.46& 70.09\\
& ViT-B-32& & 151& 90.92& 80.08& 69.20& 61.38& 63.96& 77.21& 68.79\\
% & ViT-B-32-negclip& & 151& 92.68& 85.91& 76.46& 75.20& 75.38& 88.80& 82.80\\

& RN50x4& & 178& 92.68& 82.99& 67.57& 65.04& 63.36& 79.34& 70.09\\
& RN50x16& & 291& 93.46& 82.11& 69.20& 63.01& 65.77& 80.70& 75.87\\
& ViT-L-14& & 428& 94.07& 79.19& 65.15& 60.16& 62.31& 78.32& 71.53\\
& RN50x64& & 623& 94.49& 83.50& 70.63& 61.79& 66.67& 83.27& 73.99\\


\midrule

\multirow{6}{*}{LAION~\cite{laion5b}}   
& roberta-ViT-B-32& \multirow{4}{*}{2B}& 212& 92.86& 84.90& 72.40& 63.01& 71.02& 87.34& 79.91\\
& ViT-H-14&  & 986& 96.49& 84.77& 71.76& 67.48& 73.12& 92.05& 85.84\\
& ViT-g-14& & 1367& 95.76& 85.03& 72.40& 63.01& 71.17& 91.51& 82.08\\
& ViT-bigG-14& & 2540& 96.67& 88.07& 74.75& 62.20& 74.92& 92.19& 84.54\\ \cline{3-3}
& xlm-roberta-base-ViT-B-32& \multirow{2}{*}{5B}& 366& 93.16& 84.01& 69.20& 63.41& 67.57& 87.78& 81.07\\
& xlm-roberta-large-ViT-H-14& & 1193& 96.85& 86.04& 72.05& 63.82& 72.07& 93.11& 86.13\\


\midrule

\multirow{4}{*}{DataComp~\cite{gadre2023datacomp}}   
& \texttt{small:}ViT-B-32& 13M& 151& 56.90& 56.85& 51.99& 50.81& 50.00& 53.93& 60.55\\
& \texttt{medium:}ViT-B-32& 128M& 151 & 77.00& 69.54& 57.68& 57.72& 57.06& 66.73& 64.88\\
& \texttt{large:}ViT-B-16& 1B& 150& 92.68& 79.82& 63.94& 56.10& 57.66& 84.34& 78.61\\
& \texttt{xlarge:}ViT-L-14& 13B& 428& 95.52& 84.52& 69.99& 65.04& 66.82& 91.03& 84.97\\

\bottomrule 
\end{tabular}
}
  
  \label{tab:eval_ours}

\end{table*}
We present four key findings in our evaluation over $17$ pretrained CLIP models on \method, with results reported in Table~\ref{tab:eval_ours} and visualized in Figure~\ref{fig:imagenet_correlation}.

\textbf{The best pretrained CLIP models demonstrate some compositional understanding but still have overall large rooms for improvements.} Table~\ref{tab:eval_ours} shows that the largest pretrained CLIP models, \eg, OpenAI's RN50x64, LAION's xlm-roberta-large-ViT-H-14, and DataComp's ViT-L-14, achieve near-human performance on \replaceobj. However, on \replaceobj, smaller models pretrained on small datasets still suffer from big drops in performance --- 23\% and 43\% respectively for DataComp's small and medium models --- compared to humans. Additionally, on nearly all other hard negative types, there are clear gaps (larger than 10\%) between the best model performances and human performances, showing an overall large room for improvements in current models' compositionality.

\textbf{All models struggle at identifying \swap\ hard negatives, regardless of their pertaining dataset and model size.} Among the three types of hard negatives, \swap\ hard negatives present the biggest challenge to the pretrained CLIP models, even though humans can easily tell them apart from the positive captions.  We observe in Table~\ref{tab:eval_ours} that all models demonstrate low performance on both \swapobj\ and \swapatt\ hard negatives regardless of their pretraining dataset and model sizes, with the difference from human performance reaching from 27\% to 50\%.

\textbf{Existing models are object-centric, struggling to compose attributes and relations.} We find that existing pretrained models are a lot better at composing objects than attributes or relations (Table~\ref{tab:eval_ours}). This finding holds for both \replace\ and \add\ hard negatives but not the most difficult \swap\ negatives, where models perform equally poorly on both \swapobj\ and \swapatt.
On \replace\ hard negatives, even though most models achieve human-level performance on \replaceobj, they all suffer from a drop in performance on \replaceatt\ and \replacerel, where the drop is as large as 15\% and 29\% respectively.
Similarly, on \add\ hard negatives, all models except for DataComp's small:ViT-B-32 experience a decrease in performance from \addobj\ to \addatt, with the largest difference reaching 10\%. 
% Interestingly, human performance also decreases when the perturbed atom types are attributes or relations instead of objects, but the decrease is smaller: 11\% for \add\ and only 2-3\% for \swap\ and \replace\ hard negatives.

\textbf{Models' performance on \method correlates with their ImageNet zero-shot accuracy.}
\begin{figure}[t]
    \centering
    \includegraphics[width=\linewidth]{imgs/imagenet_zeroshot_correlation.png}
    \caption{We plot pretrained vision-language models' zero-shot top-1 accuracy on ImageNet versus their retrieval recall@1 on \method, where $r$ is the Pearson correlation coefficient. This plot suggests that models' ImageNet zero-shot accuracy positively correlates with their compositionality.}
    \label{fig:imagenet_correlation}
\end{figure}
We show in Figure~\ref{fig:imagenet_correlation} that there is a positive correlation between models' performance on \method and their zero-shot accuracy on ImageNet. This correlation is moderate on \swapobj\ and \addatt\ (Pearson correlation coefficient $r = 0.78$ and $r = 0.75$ respectively) and strong on all other hard negatives ($r > 0.8$).
% This suggests that models' higher accuracy on ImageNet largely implies better compositionality as evaluated on \method.



